\begin{abstract}
Robotic manipulation of highly deformable garments presents significant challenges due to their high dimensionality and complex dynamics. This paper presents an exploratory study of diverse imitation learning strategies within the context of the ICRA 2026 LeHome Challenge. We investigate a progression of architectures and state-of-the-art visuomotor policies, evolving from Diffusion Policy through a sequence of Vision-Language-Action (VLA) models—including SmolVLA, $\pi_{0.5}$, and X-VLA—ultimately converging on Action Chunking with Transformers (ACT). Our study systematically evaluates the impact of mid-level geometric priors (colorized depth) provided as an additional input alongside RGB, robust data augmentations, and Parameter-Efficient Fine-Tuning (PEFT) using LoRA. We further analyze the effects of critical hyperparameters such as action chunk size and execution steps. Our results demonstrate that while VLA-based models show strong semantic understanding, the fine-grained bimanual coordination required for garment folding is best captured by ACT with geometric and temporal optimizations. We report per-garment success rates, showing significant performance gains across categories, with a final macro-average success rate of 61.25\% and enhanced generalization to unseen garment styles.
\end{abstract}
